%%%%%%%%%%%%%%%%%%%%%%%%%%%%%%%%%%%%%%%%%%%%%%%%%%%%%%%%%%%%%%%%%%%%%%
% Lớp 11 | Hàm số mũ và hàm số lôgarít | Phương trình, bất phương trình mũ và lôgarít
% Dạng 1D6?4-1: Điều kiện có nghiệm
% Số câu: 7  (sinh tự động bởi tach_cau.py; từ nay thầy sửa trực tiếp file này)
%%%%%%%%%%%%%%%%%%%%%%%%%%%%%%%%%%%%%%%%%%%%%%%%%%%%%%%%%%%%%%%%%%%%%%

%%%%%%%%%%%%%%%% TRẮC NGHIỆM 4 PHƯƠNG ÁN %%%%%%%%%%%%%%%%

% nguồn: Ok_Toan11_TracNghiem_Lan2
\begin{ex}%[1D6N4-1]%[ma:f0a3a7ac]
	Phương trình nào sau đây  vô nghiệm?
	 \choice
		{$2^x=1$} 
		 {$2^x=3$}
		{\True  $2^x=0$}
		{$2^x=3^x$} 
	\loigiai{ 
		Ta có  hàm số mũ luôn dương.\\
		 Vậy $2^x=0$ vô nghiệm.
	 }
\end{ex}

% nguồn: Ok_Toan11_TracNghiem_Lan2
\begin{ex}%[1D6N4-1]%[ma:0b361e02]
	Tìm tất cả các giá trị thực của $m$ để phương trình $3^x=m$ có  nghiệm thực.
	\choice 
		 {\True $m>0$}
		{$m\geq 1$} 
		{$m\geq 0$} 
		{$m\neq 0$} 
	\loigiai{ 
		Để phương trình $3^x=m$ có nghiệm thực thì  $m>0$.
	 }
\end{ex}

% nguồn: Ok_Toan11_TracNghiem_Lan2
\begin{ex}%[1D6H4-1]%[ma:328963f2]
	Tìm tất cả các giá trị của tham số m để  phương trình $\left(\dfrac{1}{3}\right)^x=m-1$ có nghiệm
	 \choice
		 {$m<0$ hoặc $m=1$}
		 {\True $m>1$}
		{$m\geq  0$}
		{$0<m\leq 1$} 
	\loigiai{ 
		Phương trình có nghiệm $\Leftrightarrow m-1>0  \Leftrightarrow m>1$.
	 }
\end{ex}

% nguồn: Ok_Toan11_TracNghiem_Lan2
\begin{ex}%[1D6H4-1]%[ma:ffe2c700]
	Có bao  nhiêu giá trị nguyên của tham số $m$ để phương trình $5^x=9-m^2$ có nghiệm thực?
	\choice 
		{$6$} 
		{\True  $5$}
		{$4$} 
		{$7$} 
	\loigiai{ 
		YCBT  $\Leftrightarrow 9-m^2 > 0 \Leftrightarrow -3<m<3$. \\
		Do $m \in \mathbb{Z}$ nên  $m \in \{-2;-1;0;1;2\}$.
	 }
\end{ex}

% nguồn: Ok_TracNghiemLop11_Lan1
\begin{ex}%[1D6H4-1]%[ma:54d60d1b]
	Tập hợp tất cả các giá trị của tham số $m$ để phương trình $\left(\dfrac{1}{3}\right)^x=m-1$ có nghiệm là 
	 \choice
		{$  (0;1]$}
		 {$(-\infty; 0) \cup \{1 \}$}
		{$[0; +\infty)$} 
		{\True  $(1; +\infty)$}
	 \loigiai{
		Phương trình $\left(\dfrac{1}{3}\right)^x=m-1$ có nghiệm $\Leftrightarrow m - 1 > 0 \Leftrightarrow  m>1$.
	} 
\end{ex}

% nguồn: Ok_Toan11_TracNghiem_Lan2
\begin{ex}%[1D6V4-1]%[ma:7e218c07]
	Tất cả các giá trị của tham số $m$ để  phương trình $3^{x^2+1}=m-1$ có nghiệm là
	 \choice
		{\True  $m\geq 4$}
		 {$m>4$}
		 {$m>1$}
		{$m\geq  1$}
	 \loigiai{
		Ta có $3^{x^2} \ge 3^0 \Leftrightarrow 3^{x^2+1} \ge 3$. \\ 
		Phương trình $3^{x^2+1}=m-1$  có nghiệm khi và chỉ khi $m-1 \ge 3 \Leftrightarrow m \ge 4$.
	 }
\end{ex}


%%%%%%%%%%%%%%%% TỰ LUẬN %%%%%%%%%%%%%%%%

% nguồn: Ok_TuanLuan_Toan11_Lan2
\begin{ex}%[1D6V4-1]%[ma:42215ce3]
	Tìm $m$ để phương trình $m\cdot2^{x^2-5x+6}+2^{1-x^2}=2\cdot2^{6-5x}+m\; \;  (1)$ có $4$ nghiệm phân biệt.
	\loigiai{ 
		Ta có 
		\begin{eqnarray*}
			m\cdot 2^{x^2-5x+6}+2^{1-x^2}=2^{7-5x}+m &\Leftrightarrow & m\cdot 2^{x^2-5x+6}+2^{1-x^2}=2^{\left(x^2-5x+6\right)+\left(1-x^2\right)}+m\\ 
			 &\Leftrightarrow & m\cdot2^{x^2-5x+6}+2^{1-x^2}=2^{x^2-5x+6}\cdot 2^{1-x^2}+m\\
			&\Leftrightarrow &  \left(2^{x^2-5x+6}-1\right)\left(2^{1-x^2}-m\right)=0\\
			&\Leftrightarrow & \hoac{&2^{x^2-5x+6}=1 \\  &2^{1-x^2}=m}\\
			&\Leftrightarrow  & \hoac{&x=3 \\&x=2 \\&2^{1-x^2}=m. \qquad (*)}
		\end{eqnarray*}
		Phương trình $(*)\Leftrightarrow \heva{&m > 0\\&1-x^2=\log_2m} \Leftrightarrow \heva{&m > 0\\ &x^2=1-\log_2m.}$\\ 
		$(1)$ có $4$ nghiệm phân biệt $\Leftrightarrow (*)$ có $2$ nghiệm phân biệt khác $2$ và  $3$.
		\[\Leftrightarrow  \heva{&{m > 0} \\
		&{1-\log_2m > 0}  \\
		 &{1-\log_2m\ne 4} \\
		 &{1-\log_2m\ne 9}
		} \Leftrightarrow  \heva{&{m > 0} \\
		&{m < 2} \\ 
		 &{m\ne \dfrac{1}{8}} \\
		&{m\ne \dfrac{1}{256}} 
		} \Leftrightarrow m\in \left(0;2\right)\setminus \left\{\dfrac{1}{8};\dfrac{1}{256} \right\}. 
		 \]
	} 
\end{ex}
